Diese Diplomarbeit dokumentiert die Weiterentwicklung, Neukonstruktion und Systemintegration des autonomen Modellfahrzeugs \emph{Tumbler~v3.0} im Rahmen des Projekts \emph{Self-Driving Robot Experiments} an der HTL BULME Graz-Gösting. Aufbauend auf den Ergebnissen einer vorangegangenen Diplomarbeit wurden zentrale mechanische, schaltungstechnische und softwareseitige Einschränkungen analysiert und durch eine durchgängig neu entwickelte Plattform ersetzt.

Im Hardwarebereich entstand eine zweistufige Elektronikarchitektur rund um den Mikrocontroller \emph{ESP32-S3-WROOM-1}. Für die Spurerkennung wurde eine vollständig neu geroutete, sechskanalige Liniensensor-Platine (\qtyproduct{76 x 32}{\milli\meter}) mit optischen Reflexkopplern (\emph{CNY70}), zwei diskreten Transistor-Stromsenken (\emph{BC817-40}) und RC-Tiefpassfiltern (\qty{1,93}{\kilo\hertz}) entwickelt und fehlerfrei geprüft (0~ERC-/0~DRC-Meldungen). Parallel dazu wurde das Fahrgestell in zehn iterativen CAD-Stufen bis zur druckfertigen Revision~v10.0 überarbeitet. Das monolithische Unibody-Chassis (\qty{220,8}{\centi\meter\cubed}, 100\,\% wasserdichtes STL-Netz für den Bambu Lab P1S) integriert acht \emph{VL53L0/1XV2}-LiDAR-Sensoren mit kollisionsfreien \qty{25}{\degree}-Sichtkegeln, einen nach oben offenen C-Führungskanal zum Einschieben des verlöteten Hauptschalters, sechzehn M2-Teardrop-Sacklöcher sowie vier verdrehsichere M3-Sechskant-Mutterntaschen.

Auf Seiten der Sensorik und Software wurden eine sequenzielle XSHUT-Initialisierung zur dynamischen I\textsuperscript{2}C-Adressvergabe (\texttt{0x30} bis \texttt{0x37}), ein mehrstufiger Median-/EMA-Filter, ein geometrischer Algorithmus zur Längsparklücken-Vermessung sowie eine deterministische FreeRTOS-Softwarearchitektur mit kaskadierter diskreter PID-Regelung (Clamping Anti-Windup) und einem zustandsbasierten Einparkautomaten konzipiert und in C++ (ESP-IDF) implementiert. Die geometrische, elektrische und kinematische Verifikation bestätigt die Funktionsfähigkeit der Gesamtarchitektur; abschließende Hardware-Labormessungen am fertig bestückten Gesamtfahrzeug sind im Testplan strukturiert vorbereitet.
